\thispagestyle{plain}
\nurvollversion{\vspace*{\fill}}
\nurplainversion{\chapter{Marts Aufzeichnungen zum Erbe des Wunderkindes}}

\textit{Jahrhunderte vor den Legenden von Andor schufen der erfinderische Zwerg Kreatok und der tempramentvolle Drache Nehal gemeinsam die vier mächtigen Schilde aus uralter Zeit, nach denen wir Schildzwerge uns noch heute benennen. Doch das Wunderkind Kreatok kreierte auch zahlreiche andere Wunderwerke: Fein ausgedachte Fallen und Verteidigungsanlagen, mit magischen Runenfolgen verzierte kunstvolle Reliefs sowie die Sphäre, ein in den Chroniken unseres Volkes fast völlig verschollen gegangenes Artefakt, welches es ermöglicht, Portale zwischen Vergangenheit und Zukunft zu spannen.}
\bigskip

\textit{In den Wirren des Unterirdischen Kriegs zwischen Schildzwergen, Drachen, Krahdern und Trollen gingen die meisten Werke Kreatoks und Nehals verloren. Erst in den Siebzigerjahren nach andorischer Zeitrechnung wurden die vier mächtigen Schilde wieder vereint, nur um bei der legendären Zerstörung der Krahder-Feste Borghorn unwiederbringlich im Lavameer zu versinken. Der Zufall – oder das Schicksal – will es, dass Kreatoks Sphäre wenige Jahre zuvor nach Jahrhunderten des verstaubten Herumliegens wieder aus der geheimen Fürstenkammer entnommen und ebenso unwiederbringlich zerschmettert wurde. Zudem wurde Kreatoks versiegelte Kultstätte kürzlich von Drachenkultisten entweiht. Ein Wunder, dass der dort gelagerte Rubin von Kreatok wieder zurückerobert werden konnte, ehe auch er vernichtet worden wäre.}
\bigskip

\textit{Wir leben in einer Ära, in welcher sich das Erbe des Wunderkindes augenscheinlich abrupt seinem Ende nähert. Dies ist an sich weder gut noch schlecht. Früher oder später findet alles ein Ende, so will es der Lauf der Dinge, und Kreatoks Kreationen haben in ihrer langen Existenz vergleichbar viel Leid verursacht wie Freude geschenkt. Doch als jemand, dem die seltene Ehre zuteilwurde, den zeitreisenden Kreatok persönlich anzutreffen, fände ich es zu schade, wenn das Wissen über die Werke des ersten Schildzwergs ebenso rasch dem Vergessen anheimfallen würde. Denn jene spielten in den legendären Geschehnissen der letzten Jahrzehnte eine erheblich größere Rolle, als selbst die sonst so allwissend scheinenden Bewahrer vom Baum der Lieder ihnen zuschreiben würden.}
\bigskip

– Aufzeichnungen von Mart, stellvertretender Fürst über Cavern in Absenz von Fürstin Marun und Fürst Kram, verfasst im Frühjahr 76 a.Z.

\nurvollversion{\vspace*{\fill}}
\newpage
\thispagestyle{plain}
\nurplainversion{\chapter{Vorwort}}

\textbf{Liebe Andori} 
\bigskip

\textbf{Die Legenden von Andor sind ein multimediales Franchise mit einer einzigartigen Fan-Community.} 
Ich kann mir kaum vorstellen, dass jemand über dieses Dokument stolpert, ohne bereits einiges zu Andor zu wissen. Falls dennoch: Die Legenden von Andor sind eine Reihe von kooperativen Brettspielen mit zahlreichen Spin-Offs sowie einer aktiven, offenen und kreativen Fan-Gemeinschaft, von der ich geehrt bin, mich dazu zählen zu dürfen.
\bigskip

\textbf{Hier herrscht Spoiler-Gefahr für die offizielle Andor-Geschichte.} 
Anspielungen sind das Brot und Butter der bärigen Geschichtchen. Nebst Spoilergefahr für alle möglichen Andor-Produkte besteht hier auch Verwirrungsgefahr für Andor-Neulinge. Für einen Überblick über die offiziellen Andorversen empfehle ich meine Chronik der Andorversen (\url{https://legenden-von-andor.de/forum/viewtopic.php?f=11&t=10314}), nachdem ihr euren ersten Kontakt mit Andor hinter euch habt.
\bigskip

\textbf{Diese Geschichtensammlung umfasst den Großteil meiner andorischen Fanfiction.} 
Die meisten dieser Geschichten spielen in einer gemeinsamen Zeitlinie, dem "bärig gebutterten Andorversum". Sie bauen aufeinander auf, sollten jedoch in beliebiger Reihenfolge gelesen werden können. Nicht dazu passende Kurzgeschichten sind am Ende dieser Sammlung angehängt.
Die Geschichten wurden ohne generative KI geschrieben. 
Sie könnten nicht existieren ohne die brillanten Werke von Michael Menzel und Stefanie Schmitt oder ohne die Tavernengemeinschaft, insbesondere die Schlafende Katze, die mich zu meinem ersten Andor-Märchen anstiftete, oder TroII und seinen "Der Ewige Rat", die umfangreichste Andor-Geschichte überhaupt, zu deren Lektüre (\url{https://legenden-von-andor.de/forum/viewtopic.php?f=26&t=11465}) ich von ganzem Herzen raten kann. Ein riesiges DANKE an euch alle!
\bigskip

\textbf{Inhaltswarnung:} 
Dies sind fantastische Erzählungen, die sich mehr um Zeitschleifen und Andor-Lore-Anspielungen kümmern als um Gewaltdarstellung oder Liebesszenen. Nichtsdestotrotz kommt es vor, dass Figuren gravierende Verletzungen erleiden und erteilen, dass sie töten und sterben, dass sie Schimpfwörter fluchen (aber nicht Schlimmeres als "Scheiße") und intime Begegnungen andeuten (aber nichts Expliziteres als Küsse). Ich tu mich schwer mit Altersempfehlungen. Kennt ihr nur die Andor-Junior-Romane (offiziell 8+), sind manche der hiesigen Geschichten vielleicht zu düster. Kamt ihr mit "Das Lied des Königs" (offiziell 14+) klar, kommt ihr wahrscheinlich auch hiermit klar. Kamt ihr mit "Varkurs Erwachen" (offiziell 16+) oder TroIIs "Der Ewige Rat" klar, dann sogar höchstwahrscheinlich.
\bigskip

\textbf{Die Geschichten halten sich nahe am offiziellen Lore und können überarbeitet werden.}
Das bärig gebutterte Andorversum hält sich sehr nahe am offiziellen Kanon, entspricht ihm jedoch nicht immer zu 100 Prozent. Das wäre auch eine unrealistische Forderung. Nicht einmal offizielle Werke bleiben völlig widerspruchsfrei zueinander und neue offizielle Werke können jederzeit den Kanon so ergänzen, dass ältere Fan-Projekte nicht mehr dazu passen. Ich behalte mir vor, in Zukunft Details dieser Geschichten zu ändern. Grammatikfehler kann ich ganz einfach korrigieren, gröbere Story-Probleme eher nicht, was euch aber ja nicht davon abhalten soll, diese zu bemängeln. Ich kann nur davon profitieren.
\bigskip

\textbf{Gebt Rückmeldung!} 
Ich schreibe diese Geschichten vor allem, weil es mir Spaß bereitet, und nicht, weil ich gut darin wäre. Ihr sollt euch auf keinen Fall durch sie hindurchquälen, wenn sie euch nicht ansprechen. Aber falls jemand Freude daran findet, könnt ihr aktiv werden (idealerweise auf \url{https://legenden-von-andor.de/forum/viewtopic.php?f=11&t=10516}): Sagt mir, was ich verbessern kann, und sagt mir, was ihr toll findet und ich bei Überarbeitungen und zukünftigen Geschichten beibehalten soll.
Feedback motiviert ungemein!
\bigskip

\textbf{Und nun sucht euch einen spannenden Titel und taucht ein in die Welt der bärig gebutterten Geschichtchen. Viel Spaß beim Schmökern, und möget ihr zum Schreiben eigener Storytexte angeregt werden.}
\bigskip


Liebe Grüße

Butterbrotbär